\chapter{Background}
\label{ch:background}

% For AI-assisted explanatory figures in this chapter, use \aifigurecaption{...}
% so the standardized disclosure is included directly in the caption.

This chapter introduces autoregressive image modeling and vector quantization,
then describes early and recent image tokenizers.

\section{Autoregressive Models: From Text to Images}
\label{sec:ar_models_text_to_images}

An autoregressive model represents a joint distribution by predicting each
element from those that precede it. A token is a discrete text unit, such as a
word or subword, drawn from a finite vocabulary. Let
$s=(s_1,\ldots,s_T)$ be a sequence of tokens, where $s_i$ is the token at
position $i$ and $T$ is the sequence length. The chain rule gives
\begin{equation}
    p(s)=\prod_{i=1}^{T}p(s_i\mid s_{<i}),
    \qquad s_{<i}=(s_1,\ldots,s_{i-1}).
\end{equation}
This factorization has long been used for statistical language
modeling~\cite{bengio2003neural}. A text tokenizer converts text into a sequence
of tokens drawn from a finite vocabulary, and a language model learns the
conditional distribution of the next token. Transformers made this approach
effective at scale by using self-attention to model dependencies between
tokens~\cite{vaswani2017attention}. Later work found that language-model loss
follows empirical power laws as the amount of data, model size, and training
compute increase~\cite{kaplan2020scaling}. Large autoregressive language models
then demonstrated broad capabilities under the same training
objective~\cite{brown2020language,touvron2023llama}.

Applying the same principle to images is less direct. An image is a
two-dimensional array rather than an ordered sequence, so an autoregressive
model must first choose both a prediction unit and an ordering. The most literal
choice is to discretize pixel intensities and traverse them in raster order,
from left to right and top to bottom. Early works such as PixelRNN and PixelCNN
showed that this produces tractable image likelihoods and coherent
samples~\cite{oord2016pixel}. Image Transformer later applied self-attention to
autoregressive pixel generation while restricting attention to local
neighborhoods to control its cost~\cite{parmar2018image}.

Pixel-level modeling is conceptually simple but computationally demanding. A
$256\times256$ RGB image contains $196{,}608$ channel values, and generation
must resolve low-level color statistics as well as long-range structure.
Moreover, full self-attention on a sequence of length $T$ has
$\mathcal{O}(T^2)$ time and memory cost. A learned image tokenizer addresses
both problems by mapping the image to a shorter grid of discrete symbols. The
token grids studied in this thesis are flattened in raster order. Tokenization
reduces the number of prediction units but does not choose a different
ordering. Figure~\ref{fig:pixel_vs_token_sequence} contrasts the resulting
sequence lengths.

\begin{figure}[t]
    \centering
    \resizebox{0.94\textwidth}{!}{\begin{tikzpicture}[
    font=\sffamily\small,
    panel/.style={draw=black!35, rounded corners=2pt, minimum width=70mm,
                  minimum height=34mm, align=center, inner sep=6pt},
    sequence/.style={draw=black!55, minimum height=5mm, minimum width=7mm,
                     inner sep=0pt},
    arrow/.style={-{Latex[length=2mm]}, semithick, black!65}
]
\node[panel, fill=red!4] (pixels) at (-38mm,0) {};
\node[font=\bfseries, anchor=north] at ([yshift=-3mm]pixels.north)
    {Pixel-level sequence};
\node[align=center] (pixelshape) at ([yshift=3mm]pixels.center)
    {$x\in[0,1]^{256\times256\times3}$\\[1mm]
     $T_{\mathrm{pixel}}=256\cdot256\cdot3=196{,}608$ scalar values};
\node[anchor=south, align=center, font=\scriptsize] at ([yshift=2mm]pixels.south)
    {Raster order over all RGB channel values};

\node[panel, fill=blue!4] (tokens) at (38mm,0) {};
\node[font=\bfseries, anchor=north] at ([yshift=-3mm]tokens.north)
    {Token-level sequence};
\node[align=center] (tokenshape) at ([yshift=3mm]tokens.center)
    {$k\in\{1,\ldots,K\}^{16\times16}$\\[1mm]
     $T_{\mathrm{token}}=16\cdot16=256$ discrete tokens};
\node[anchor=south, align=center, font=\scriptsize] at ([yshift=2mm]tokens.south)
    {Raster order over the learned token grid};

\draw[arrow] ([xshift=-9mm,yshift=-7mm]pixels.center) --
    node[below=1mm, font=\scriptsize] {$196{,}608$ positions}
    ([xshift=9mm,yshift=-7mm]pixels.center);
\draw[arrow] ([xshift=-9mm,yshift=-7mm]tokens.center) --
    node[below=1mm, font=\scriptsize] {$256$ positions}
    ([xshift=9mm,yshift=-7mm]tokens.center);

\node[font=\sffamily\bfseries] at (0,-25mm)
    {Sequence-length reduction: $196{,}608/256=768$};
\end{tikzpicture}
}
    \caption{Pixel-level flattening produces much longer autoregressive sequences than token-level flattening. A learned tokenizer replaces groups of pixel values with discrete visual symbols.}
    \label{fig:pixel_vs_token_sequence}
\end{figure}

Autoregressive image generation therefore has two stages, illustrated in
Figure~\ref{fig:two_stage_autoregressive_pipeline}. A tokenizer first converts
images into compact token grids and reconstructs images from them. An
autoregressive model is then trained on the resulting token sequences. This
second-stage model is also called a \emph{prior}: it models the distribution of
token sequences and can sample a sequence for the decoder.

\begin{figure}[t]
    \centering
    \resizebox{0.96\textwidth}{!}{\begin{tikzpicture}[
    font=\sffamily\scriptsize,
    block/.style={draw=black!50, rounded corners=2pt, minimum width=22mm,
                  minimum height=9mm, align=center, fill=blue!7},
    data/.style={draw=black!45, minimum width=18mm, minimum height=9mm,
                 align=center, fill=black!2},
    prior/.style={block, fill=orange!10, minimum width=29mm},
    arrow/.style={-{Latex[length=2mm]}, semithick, black!65}
]
\node[font=\bfseries\small, anchor=west] (s1) at (-75mm,15mm)
    {Stage 1: learn and then freeze the tokenizer};
\node[data, anchor=west] (x) at (-75mm,3mm) {image\\$x$};
\node[block, right=5mm of x] (enc) {encoder\\$E$};
\node[data, right=5mm of enc] (grid) {token grid\\$k\in\{1,\ldots,K\}^{16\times16}$};
\node[block, right=5mm of grid] (dec) {decoder\\$D$};
\node[data, right=5mm of dec] (recon) {reconstruction\\$\hat{x}$};
\draw[arrow] (x) -- (enc);
\draw[arrow] (enc) -- (grid);
\draw[arrow] (grid) -- (dec);
\draw[arrow] (dec) -- (recon);

\node[font=\bfseries\small, anchor=west] (s2) at (-75mm,-18mm)
    {Stage 2: learn an autoregressive prior over tokenizer outputs};
\node[data, anchor=west] (trainseq) at (-75mm,-30mm)
    {training sequence\\$s=(k_1,\ldots,k_{256})$};
\node[prior, right=5mm of trainseq] (ar)
    {autoregressive prior\\$p(s)=\prod_i p(k_i\mid k_{<i})$};
\node[data, right=5mm of ar] (sample)
    {sampled sequence\\$\tilde{s}$};
\node[block, right=5mm of sample] (dec2) {frozen\\decoder $D$};
\node[data, right=5mm of dec2] (generated) {generated image\\$\tilde{x}$};
\draw[arrow] (trainseq) -- (ar);
\draw[arrow] (ar) -- (sample);
\draw[arrow] (sample) -- (dec2);
\draw[arrow] (dec2) -- (generated);
\draw[arrow, dashed] (grid.south) -- ++(0,-5mm) -| (trainseq.north);
\end{tikzpicture}
}
    \caption{Two-stage autoregressive image generation. The tokenizer maps an image to a compact token grid; an autoregressive prior models the flattened token sequence; the tokenizer decoder maps sampled tokens back to image space.}
    \label{fig:two_stage_autoregressive_pipeline}
\end{figure}

Image tokenizers may use convolutions, attention, or both. Their function is to
represent images with discrete tokens and reconstruct images from them. The
autoregressive model has a different function: it predicts those tokens in
sequence. The tokenizer determines which image information reaches that model,
the sequence length, and the vocabulary used for next-token prediction.

\section{Image Tokens and Vector Quantization}
\label{sec:image_tokens_vq}

Vector quantization replaces continuous feature vectors with entries from a
finite codebook. We introduce this operation using VQ-VAE, which learns
discrete latent representations by assigning encoder outputs to a learned
codebook~\cite{oord2017neural}. A separate autoregressive model can then model
the resulting token sequence. VQ-VAE-2 later showed that hierarchical discrete
latents and stronger priors could support high-fidelity large-scale image
generation~\cite{razavi2019vqvae2}.

VQ-VAE contains an encoder, a vector quantizer, and a decoder. Given an input
image $x\in\mathbb{R}^{H\times W\times C}$, the encoder produces a continuous
feature grid $z_e(x)\in\mathbb{R}^{h\times w\times d}$. The full path is shown
in Figure~\ref{fig:vqvae_pipeline}.

\begin{figure}[t]
    \centering
    \resizebox{0.9\textwidth}{!}{\begin{tikzpicture}[
    font=\sffamily\small,
    data/.style={draw=black!45, rounded corners=2pt, fill=black!2,
                 minimum width=24mm, minimum height=14mm, align=center},
    block/.style={draw=black!50, rounded corners=2pt, fill=blue!7,
                  minimum width=27mm, minimum height=14mm, align=center},
    arrow/.style={-{Latex[length=2mm]}, semithick, black!65}
]
\node[data] (input) {input image\\$x$};
\node[block, right=7mm of input] (encoder) {encoder\\$E$};
\node[data, right=7mm of encoder] (features) {continuous features\\$z_e(x)$};
\node[block, right=7mm of features] (quantizer) {vector quantizer\\$Q$};
\node[data, right=7mm of quantizer] (tokens) {tokens $k$\\vectors $z_q(x)$};
\node[block, right=7mm of tokens] (decoder) {decoder\\$D$};
\node[data, right=7mm of decoder] (output) {reconstruction\\$\hat{x}$};
\draw[arrow] (input) -- (encoder);
\draw[arrow] (encoder) -- (features);
\draw[arrow] (features) -- (quantizer);
\draw[arrow] (quantizer) -- (tokens);
\draw[arrow] (tokens) -- (decoder);
\draw[arrow] (decoder) -- (output);
\end{tikzpicture}
}
    \aifigurecaption{VQ-VAE representation path. The encoder produces continuous
    features, the vector quantizer returns integer tokens and their associated
    codebook vectors, and the decoder reconstructs the image. Based on van den
    Oord et al.~\cite{oord2017neural}.}
    \label{fig:vqvae_pipeline}
\end{figure}

At each spatial position $(u,v)$, the vector $z_e(x)_{u,v}\in\mathbb{R}^{d}$ is assigned to its nearest entry in a learned codebook
$\mathcal{E}=\{e_1,\ldots,e_K\}$, where $e_j\in\mathbb{R}^{d}$. Quantization produces an integer token grid $k\in\{1,\ldots,K\}^{h\times w}$ and a quantized latent tensor $z_q(x)\in\mathbb{R}^{h\times w\times d}$:
\begin{align}
    z_e(x) &= E(x), \\
    k_{u,v} &= \underset{j\in\{1,\ldots,K\}}{\arg\min}\,
        \lVert z_e(x)_{u,v}-e_j\rVert_2, \\
    z_q(x)_{u,v} &= e_{k_{u,v}}, \\
    \hat{x} &= D\!\left(z_q(x)\right).
\end{align}
Here, $D$ is the decoder and $\hat{x}$ is the reconstruction. Figure~\ref{fig:vq_clustering_codebook} gives the corresponding geometric interpretation: codebook entries act as learned representative points in latent space.

\begin{figure}[t]
    \centering
    \includegraphics[width=0.94\textwidth]{ch02_fig_2_2a_vq_clustering_codebook.png}
    \aifigurecaption{Vector quantization as clustering in latent space. Continuous encoder outputs are assigned to their nearest learned codebook entry, and the entry index becomes a discrete token.}
    \label{fig:vq_clustering_codebook}
\end{figure}

The decoder receives $z_q(x)$ rather than the continuous encoder output. Image information must therefore pass through a finite vocabulary of $K$ vectors. The same representation has two roles: the quantized vectors are the decoder input, while their integer indices form the discrete sequence modeled by the autoregressive prior. Figure~\ref{fig:feature_map_to_token_sequence} summarizes this conversion.

\begin{figure}[t]
    \centering
    \resizebox{0.94\textwidth}{!}{\begin{tikzpicture}[
    font=\sffamily\scriptsize,
    block/.style={draw=black!50, rounded corners=2pt, minimum width=25mm,
                  minimum height=14mm, align=center, fill=blue!7},
    data/.style={draw=black!45, minimum width=27mm, minimum height=14mm,
                 align=center, fill=black!2},
    quant/.style={block, fill=red!7, minimum width=31mm},
    arrow/.style={-{Latex[length=2mm]}, semithick, black!65}
]
\node[data] (image) at (-66mm,0) {image\\$x\in\mathbb{R}^{H\times W\times C}$};
\node[block, right=5mm of image] (encoder) {encoder\\$E$};
\node[data, right=5mm of encoder] (latent)
    {continuous grid\\$z_e(x)\in\mathbb{R}^{h\times w\times d}$};
\node[quant, right=5mm of latent] (nearest)
    {nearest codebook entry\\$k_{u,v}=\arg\min_j\|z_{e,u,v}-e_j\|_2$};
\node[data, right=5mm of nearest] (indices)
    {index grid\\$k\in\{1,\ldots,K\}^{h\times w}$};
\node[data, below=12mm of indices] (sequence)
    {raster sequence\\$s=(k_{1,1},k_{1,2},\ldots,k_{h,w})$};
\draw[arrow] (image) -- (encoder);
\draw[arrow] (encoder) -- (latent);
\draw[arrow] (latent) -- (nearest);
\draw[arrow] (nearest) -- (indices);
\draw[arrow] (indices) -- node[right=1mm] {flatten} (sequence);
\node[below=7mm of sequence, align=center, font=\scriptsize, text width=82mm]
    {For the evaluated checkpoints, $h=w=16$ and $K=16{,}384$: the grid has
     $256$ positions, and each token is one of $16{,}384$ codebook indices.};
\end{tikzpicture}
}
    \caption{From a continuous feature map to a discrete sequence. The encoder produces a latent grid, vector quantization replaces each latent vector with a codebook index, and the token grid is flattened for sequence modeling.}
    \label{fig:feature_map_to_token_sequence}
\end{figure}

\subsection{Learning the Discrete Representation}

VQ-VAE is trained by gradient descent, but nearest-neighbor assignment is not
differentiable. It therefore uses a straight-through estimator. During the
forward pass the decoder receives quantized vectors; during backpropagation,
the gradient with respect to the quantized representation is copied to the
encoder output~\cite{bengio2013estimating,oord2017neural}. The encoder, decoder,
and codebook are learned together.

The objective contains a reconstruction loss and two additional terms for
learning the codebook and encoder representation:
\begin{equation}
    \mathcal{L}_{\mathrm{VQ}}
    =
    \underbrace{-\log p\!\left(x\mid z_q(x)\right)}_{\mathcal{L}_{\mathrm{rec}}}
    + \underbrace{\left\lVert \operatorname{sg}[z_e(x)]-z_q(x)\right\rVert_2^2}_{\mathcal{L}_{\mathrm{cod}}}
    + \underbrace{\beta\left\lVert z_e(x)-\operatorname{sg}[z_q(x)]\right\rVert_2^2}_{\mathcal{L}_{\mathrm{comm}}}.
\end{equation}
The reconstruction term follows the likelihood convention in the original
VQ-VAE formulation~\cite{oord2017neural}. The stop-gradient operator
$\operatorname{sg}[\cdot]$ is the identity in the forward pass and has zero
derivative. The codebook term $\mathcal{L}_{\mathrm{cod}}$ moves embeddings
toward encoder outputs. The commitment term $\mathcal{L}_{\mathrm{comm}}$
keeps encoder outputs close to their assigned embeddings.

The codebook contains $K$ entries, but a dataset may activate only a subset.
Codebook collapse occurs when many entries receive few or no assignments, which
reduces the empirical vocabulary available to the autoregressive
model~\cite{lancucki2020robust,yu2021vector,mentzer2023finite}. Chapter~\ref{ch:rq1}
measures nominal and empirical vocabulary directly.

Nearest-neighbor assignment also partitions latent space into discrete regions.
A small change in an encoder output can cross a region boundary and change its
token abruptly, as illustrated in Figure~\ref{fig:vq_boundary_crossing}.

\begin{figure}[t]
    \centering
    \resizebox{0.9\textwidth}{!}{\begin{tikzpicture}[
    font=\sffamily\small,
    point/.style={circle, fill=black, inner sep=1.7pt},
    centroid/.style={circle, draw=red!70!black, fill=red!65, inner sep=2.6pt},
    move/.style={-{Latex[length=2.3mm]}, very thick, red!70!black}
]
\draw[-{Latex[length=2mm]}] (-55mm,-24mm) -- (55mm,-24mm)
    node[below] {$z_1$};
\draw[-{Latex[length=2mm]}] (-55mm,-24mm) -- (-55mm,25mm)
    node[left] {$z_2$};
\draw[dashed, semithick] (0,-22mm) -- (0,23mm)
    node[above, align=center, font=\scriptsize] {equal-distance boundary\\$\|z-e_a\|_2=\|z-e_b\|_2$};
\node[centroid, label=below left:$e_a$] (ea) at (-28mm,-1mm) {};
\node[centroid, label=below right:$e_b$] (eb) at (28mm,6mm) {};
\node[point] (before) at (-7mm,2mm) {};
\node[above left=1mm of before] {$z$};
\node[point] (after) at (8mm,5mm) {};
\node[above right=1mm of after] {$z'=z+\delta$};
\draw[move] (before) -- node[above] {$\delta$} (after);
\draw[thin, black!45] (before) -- (ea);
\draw[thin, black!45] (after) -- (eb);
\node[align=center, below=7mm of before, xshift=-15mm, font=\scriptsize]
    {$\|z-e_a\|_2<\|z-e_b\|_2$\\$Q(z)=a$};
\node[align=center, below=7mm of after, xshift=15mm, font=\scriptsize]
    {$\|z'-e_b\|_2<\|z'-e_a\|_2$\\$Q(z')=b$};
\node[draw=red!65, rounded corners=2pt, fill=red!5, align=center,
      right=13mm of eb, text width=34mm, inner sep=5pt]
    {A small continuous displacement can produce the discrete change $a\rightarrow b$.};
\end{tikzpicture}
}
    \caption{Nearest-neighbor regions create discontinuities in latent space. A small displacement of an encoder output can cross a quantization boundary and change its assigned token.}
    \label{fig:vq_boundary_crossing}
\end{figure}

\section{VQGAN and Taming Transformers}
\label{sec:vqgan_taming}

VQ-VAE showed how to learn a discrete token space that an autoregressive model
can model, with a decoder that converts sampled token sequences into
images~\cite{oord2017neural}. Its reconstruction objective, however, does not
necessarily preserve perceptually important high-frequency detail. When several
sharp reconstructions are plausible, minimizing an average pixel error can
favor a blurred result. Esser et al.\ address this limitation with a
vector-quantized generative adversarial network, commonly called VQGAN, and use
its discrete representation to make transformer-based high-resolution
synthesis practical~\cite{esser2021taming}.

VQGAN retains the encoder--quantizer--decoder structure of VQ-VAE but changes
the training objective. In the implementation used by the evaluated VQGAN
checkpoint, the generator objective has four components:
\begin{equation}
    \mathcal{L}_{\mathrm{VQGAN}}
    =
    \mathcal{L}_{\mathrm{rec}}
    + \mathcal{L}_{\mathrm{perc}}
    + \lambda\,\mathcal{L}_{\mathrm{adv}}
    + \mathcal{L}_{\mathrm{codebook}}.
\end{equation}
This expanded implementation view separates the pixel reconstruction,
perceptual, adversarial, and codebook terms; their configured and adaptive
weights are omitted here for readability. The perceptual term compares
activations of a pretrained network rather than requiring exact pixel
agreement. A patch-based discriminator supplies the adversarial term. These
objectives encourage the decoder to preserve perceptually salient structure and
synthesize plausible fine detail. Figure~\ref{fig:adversarial_and_perceptual_losses}
shows how the losses supervise the reconstructed image.

\begin{figure}[t]
    \centering
    \resizebox{0.94\textwidth}{!}{\begin{tikzpicture}[
    font=\sffamily\scriptsize,
    block/.style={draw=black!50, rounded corners=2pt, minimum width=21mm,
                  minimum height=10mm, align=center, fill=blue!7},
    data/.style={draw=black!45, minimum width=17mm, minimum height=10mm,
                 align=center, fill=black!2},
    loss/.style={draw=red!55!black, rounded corners=2pt, minimum width=29mm,
                 minimum height=9mm, align=center, fill=red!5},
    arrow/.style={-{Latex[length=2mm]}, semithick, black!65},
    supervise/.style={-{Latex[length=2mm]}, dashed, semithick, red!65!black}
]
\node[data] (x) at (-68mm,0) {real image\\$x$};
\node[block, right=5mm of x] (enc) {encoder\\$E$};
\node[block, right=5mm of enc, fill=orange!10] (quant) {quantizer\\$Q_{\mathcal E}$};
\node[block, right=5mm of quant] (dec) {decoder\\$D$};
\node[data, right=5mm of dec] (xhat) {reconstruction\\$\hat{x}$};
\draw[arrow] (x) -- (enc);
\draw[arrow] (enc) -- (quant);
\draw[arrow] (quant) -- (dec);
\draw[arrow] (dec) -- (xhat);

\node[loss, below=18mm of enc, xshift=-8mm] (rec)
    {reconstruction loss\\$\mathcal L_{\mathrm{rec}}(x,\hat{x})$};
\node[loss, below=18mm of quant, xshift=13mm] (perc)
    {perceptual loss\\$\|\phi(x)-\phi(\hat{x})\|$};
\node[block, above=17mm of xhat, fill=orange!10] (disc)
    {patch discriminator\\$C$};
\node[loss, right=5mm of disc] (adv)
    {adversarial loss\\$\mathcal L_{\mathrm{adv}}$};
\node[loss, above=17mm of quant] (codebook)
    {codebook and commitment\\$\mathcal L_{\mathrm{codebook}}$};

\draw[supervise] (x.south) -- ++(0,-7mm) -| (rec.north west);
\draw[supervise] (xhat.south) -- ++(0,-7mm) -| (perc.north east);
\node[below=2mm of rec, align=center] {inputs: $x$, $\hat{x}$};
\node[below=2mm of perc, align=center] {inputs: $x$, $\hat{x}$};
\draw[arrow] (x.north) -- ++(0,31mm) -| (disc.north);
\draw[arrow] (xhat.north) -- (disc.south);
\draw[supervise] (disc) -- (adv);
\draw[supervise] (enc.north) |- (codebook.west);
\draw[supervise] (quant.north) -- (codebook.south);

\node[below=15mm of perc, align=center, font=\sffamily\bfseries]
    {$\mathcal L_{\mathrm{VQGAN}}=\mathcal L_{\mathrm{rec}}+\mathcal L_{\mathrm{perc}}
     +\lambda\mathcal L_{\mathrm{adv}}+\mathcal L_{\mathrm{codebook}}$};
\end{tikzpicture}
}
    \caption{VQGAN training combines reconstruction, perceptual, and adversarial supervision. Reconstruction and perceptual losses compare the input and reconstruction, while a patch discriminator encourages locally realistic outputs.}
    \label{fig:adversarial_and_perceptual_losses}
\end{figure}

The first-stage model is convolutional, with non-local attention at its lowest-resolution level. This architecture combines local feature extraction with a limited mechanism for contextual interaction before quantization. Once trained, the tokenizer is frozen and images are represented by their codebook indices. A transformer is then trained autoregressively on the resulting sequences. The original architecture is reproduced in Figure~\ref{fig:vqgan_first_stage_model}.

\begin{figure}[t]
    \centering
    \includegraphics[width=0.94\textwidth]{ch02_taming_transformers_architecture.png}
    \caption{Taming Transformers architecture, comprising the VQGAN tokenizer, learned codebook, autoregressive transformer, and patch discriminator. Reproduced from Esser et al.~\cite{esser2021taming}.}
    \label{fig:vqgan_first_stage_model}
\end{figure}

A smaller token grid reduces the sequence length processed by the transformer,
but forces each token to carry more image information and can reduce
reconstruction fidelity. Esser et al.\ study this trade-off through different
codebook sizes and latent resolutions~\cite{esser2021taming}.

\section{\texorpdfstring{\mbox{LlamaGen}}{LlamaGen}}
\label{sec:llamagen_tokenizer}

After VQGAN, subsequent work improved autoregressive image generation. DALL-E,
for example, models text and image tokens as a single autoregressive
stream~\cite{ramesh2021zero}.

LlamaGen revisits this approach using a Llama-style transformer and reports
competitive results against diffusion baselines through tokenizer and model
design~\cite{sun2024llamagen}. The system preserves the two-stage design: a
vector-quantized tokenizer converts images to discrete grids, and an
autoregressive model predicts their tokens.

A large nominal codebook is of limited value if the tokenizer uses only a small
fraction of its entries. LlamaGen addresses this problem with low-dimensional
codebook embeddings and L2-normalized lookup~\cite{sun2024llamagen}. Reducing
the embedding dimension makes nearest-neighbor assignment less sparse, while
normalization places encoder projections and codebook entries on a common unit
sphere before assignment. LlamaGen also evaluates longer token sequences as a
separate design choice. This thesis evaluates its VQ-16 tokenizer.

\section{Representation Properties Beyond Reconstruction Quality}
\label{sec:reconstruction_to_behavior}

A tokenizer is usually evaluated first through reconstruction: given $x$, it produces
$\hat{x}=D(Q(E(x)))$, and the two images are compared. Reconstruction is essential because any information discarded by the tokenizer is unavailable to the generative model. It is not, however, a complete description of the induced discrete representation.

The tokenizer also determines the effective vocabulary, token frequencies, and
spatial response to perturbations. Two tokenizers with similar reconstruction
quality may use different fractions of their codebooks. A small image
perturbation may leave one token grid largely unchanged while altering many
entries in another. A local token edit may remain local or propagate through
the decoder. Chapters~\ref{ch:rq1}, \ref{ch:rq3}, and \ref{ch:rq4} measure these
properties because the autoregressive model learns the distribution created by
the tokenizer rather than the original pixel distribution directly.

VQGAN and LlamaGen represent two stages in the development of discrete image
tokenization. Both learn a finite visual vocabulary and expose token grids to
an autoregressive model, but they differ in tokenizer design and in the
empirical distributions induced over their vocabularies. Their comparison
separates behavior shared by vector-quantized tokenizers from behavior tied to
a particular design.
